\section{From projective lines to the field}
\label{sec:fields}

Proposition~\ref{prop:lines} has recovered a bijection of projective
spaces that preserves both lines and the multiplicative group law.
The fundamental theorem of projective geometry
\citep[Section~8.2, Theorem~8.2]{Hall2015} recovers addition from
the lines; the group law then forces the resulting additive map to be
multiplicative. This last passage also appears in
\citet[Section~4, Lemma~29]{CadoretPirutka2021}. We include the coordinate argument
to make both the reconstruction and its uniqueness explicit.

\begin{proposition}[Projective lifting]\label{prop:projective-lift}
Let $E$ and $E'$ be finite-dimensional vector spaces over fields $H$
and $C$, respectively, with dimensions at least three. Every bijection
$\psi:\PP_H(E)\to\PP_C(E')$ that maps lines onto lines is induced by
an additive bijection $S:E\to E'$ satisfying
\[
 S(av)=\sigma(a)S(v)\qquad(a\in H,\ v\in E)
\]
for a field isomorphism $\sigma:H\to C$. The map $\sigma$ is unique,
and $S$ is unique up to multiplication by an element of $C^\times$.
\end{proposition}

\begin{proof}
The inverse also maps lines onto lines: a target line is the image of
the line joining the preimages of any two distinct points on it.
Projective spans are obtained by repeatedly joining pairs of points
by lines, so $\psi$ and its inverse preserve spans. Minimal spanning
sets therefore correspond, and $\dim_H E=\dim_C E'=N+1$ for some
$N\ge2$. Choose a basis $e_0,\ldots,e_N$ of $E$ and representatives
$e'_0,\ldots,e'_N$ of their images. They form a basis of $E'$.
Keeping $e'_0$ fixed, scale each $e'_i$, $i\ge1$, so that
\[
 \psi([e_0+e_i])=[e'_0+e'_i].
\]
This is possible because $[e_0+e_i]$ lies on the line joining
$[e_0]$ to $[e_i]$ and differs from both endpoints.

The hyperplanes spanned by $e_1,\ldots,e_N$ and by
$e'_1,\ldots,e'_N$ correspond. On their affine complements use the
coordinates
\[
 (t_1,\ldots,t_N)\longleftrightarrow
 [e_0+t_1e_1+\cdots+t_Ne_N].
\]
The induced affine bijection preserves parallelism: parallel lines
have the same point on the corresponding hyperplane at infinity.
On the $i$th coordinate axis it is a bijection
$\sigma_i:H\to C$ fixing $0$ and $1$. Varying only the $j$th input
coordinate moves along a line with point at infinity $[e_j]$, whose
image has point at infinity $[e'_j]$. Hence all other output coordinates
are unchanged. It follows that the affine map is
\[
 (t_1,\ldots,t_N)\longmapsto
 (\sigma_1(t_1),\ldots,\sigma_N(t_N)).
\]
In each plane of two coordinate axes, the diagonal through $(0,0)$
and $(1,1)$ maps to the corresponding diagonal. Thus all $\sigma_i$ equal one bijection
$\sigma:H\to C$. In the first two coordinates, the line
$(t,t+a)$ is parallel to this diagonal and passes through $(0,a)$.
Its image therefore gives
\[
 \sigma(t+a)=\sigma(t)+\sigma(a).
\]
The line $(t,at)$ passes through $(0,0)$ and $(1,a)$, so its image gives
\[
 \sigma(at)=\sigma(a)\sigma(t).
\]
Since $\sigma(1)=1$, it is a field isomorphism.

Define
$S(\sum_{i=0}^N t_i e_i)=\sum_{i=0}^N\sigma(t_i)e'_i$.
This is an additive, $\sigma$-semilinear bijection inducing $\psi$ on
the affine chart. Every point at infinity is the direction of an
affine line through the origin, so it induces $\psi$ there as well.

For uniqueness, suppose that $T$ is another semilinear lift, possibly
with a different field isomorphism. For every nonzero $v\in E$ write
$T(v)=c_vS(v)$ with $c_v\in C^\times$. If $v,w$ are independent,
additivity applied to $v+w$ and independence of $S(v),S(w)$ give
$c_v=c_w=c_{v+w}$. Two dependent nonzero vectors can each be compared
to a vector outside their common line. Thus $c_v$ is a single scalar
$c$, and $T=cS$. Comparing the images of $av$ then shows that the two
field isomorphisms also agree.
\end{proof}

We will also recover the prescribed algebraically closed base directly
from the field.

\begin{lemma}\label{lem:perfect-core}
If $F/\kappa$ is a finitely generated extension of an algebraically
closed field of characteristic $p$, then
\[
 \bigcap_{m\ge1} F^{p^m}=\kappa.
\]
\end{lemma}

\begin{proof}
Perfectness gives the inclusion from right to left. If $f\in
F\setminus\kappa$, algebraic closedness makes $f$ transcendental over
$\kappa$. Complete $f$ to a transcendence basis
$f,t_2,\ldots,t_d$, and put
$B=\kappa(t_2,\ldots,t_d)$ and $E=B(f)$. The extension $F/E$ is
finite; write $D=[F:E]$. For every $m\ge1$, the polynomial
$T^{p^m}-f$ is Eisenstein at the prime $f$ of $B[f]$, and is therefore
irreducible over $E$. If $f=g^{p^m}$ with $g\in F$, it follows that
\[
 p^m=[E(g):E]\le D.
\]
This fails for all sufficiently large $m$, so $f$ does not belong to
the intersection.
\end{proof}

\begin{proof}[Proof of Theorem~\ref{thm:main}]
A field isomorphism $\alpha:K\to L$ sends $K^p$ to $L^p$ and
Steinberg pairs to Steinberg pairs. It therefore induces an element
of $\Isom_{\mathrm M}(V_K,V_L)$, and the canonical map in the theorem
is well defined.

For surjectivity, take $\Theta\in\Isom_{\mathrm M}(V_K,V_L)$.
Proposition~\ref{prop:lines} gives $n\in\Fp^\times$ such that
$\Psi=n\Theta$ maps lines onto lines. Put $H=K^p$ and $C=L^p$.
By Lemma~\ref{lem:p-degree}, $\dim_H K$ and $\dim_C L$ are at least
$p^2\ge4$, so
Proposition~\ref{prop:projective-lift} gives a semilinear lift
$S:K\to L$. Since $\Psi([1])=[1]$, scale $S$ to arrange $S(1)=1$.
For each $f\in K^\times$, the two semilinear bijections
\[
 w\longmapsto S(fw),\qquad w\longmapsto S(f)S(w)
\]
induce the same projective map: the equality is precisely
$\Psi([fw])=\Psi([f])+\Psi([w])$ in the additive notation for $V_L$.
Their values at $1$ agree, so the uniqueness assertion in
Proposition~\ref{prop:projective-lift} makes them equal. Thus $S$ is
multiplicative, and hence a field isomorphism. It maps every
$K^{p^m}$ onto $L^{p^m}$, so Lemma~\ref{lem:perfect-core} gives
$S(k)=l$. Its induced map is $n\Theta$, proving surjectivity onto the
scalar quotient.

For injectivity, suppose that field isomorphisms $\alpha,\beta:K\to L$
induce maps with $\beta_*=r\alpha_*$ for $r\in\Fp^\times$.
Both preserve projective lines, so multiplication by $r$ on $V_L$
also preserves lines. Lemma~\ref{lem:scalar-line} gives $r=1$.
The maps $\alpha$ and $\beta$ are therefore semilinear lifts of the
same projective bijection. Proposition~\ref{prop:projective-lift}
makes them scalar multiples, and $\alpha(1)=\beta(1)=1$ makes them
equal.
\end{proof}

\begin{remark}[Frobenius and the target quotient]
There is no further Frobenius ambiguity. For a field $F/\kappa$ as
in the theorem, the Frobenius map $F\to F$, $f\mapsto f^p$, is not
surjective and induces the zero map on $V_F$. By contrast,
$F\to F^p$ is a field isomorphism, whose induced isomorphism has target
\[
 V_{F^p}=(F^p)^\times/(F^{p^2})^\times.
\]
It is already covered by the theorem, with its actual target field
$F^p$.
\end{remark}
